% Generated draft: requires mathematical and source review.
\documentclass[11pt]{article}
\usepackage[margin=25mm]{geometry}
\usepackage{fontspec}
\setmainfont{Latin Modern Roman}
\usepackage{amsmath,amssymb,mathtools,unicode-math}
\setmathfont{Latin Modern Math}
\usepackage{enumitem,needspace,xurl,hyperref}
\setlength{\parindent}{0pt}
\setlength{\parskip}{5pt}
\setlength{\emergencystretch}{3em}
\newcommand{\N}{\mathbb{N}}
\newcommand{\Z}{\mathbb{Z}}
\newcommand{\Q}{\mathbb{Q}}
\newcommand{\R}{\mathbb{R}}
\newcommand{\C}{\mathbb{C}}
\newcommand{\F}{\mathbb{F}}
\newcommand{\A}{\mathbb{A}}
\newcommand{\PP}{\mathbb{P}}
\newcommand{\E}{\mathbb{E}}
\newcommand{\eps}{\varepsilon}
\newcommand{\dd}{\,\mathrm{d}}
\begin{document}
\section*{UnsolvedMath ID: 30001380}
\textbf{Source problem code:} catalogue entry
\subsection{Definitions and mathematical statement}
For a positive integer \(n\), let \(P\subseteq\R^n\) be a polytope, interpreted in the usual sense of a bounded convex hull of finitely many points. Define its integer-point count by
\[
N(P):=\lvert P\cap\Z^n\rvert
=\#\{(z_1,\ldots,z_n)\in P:z_i\in\Z\text{ for every }i\}.
\]
All integer points belonging to \(P\), including boundary points, are counted. No full-dimensionality or rationality assumption is supplied.

The source asks: how can one compute \(N(P)\) exactly, or estimate \(N(P)\)? This is a general methodological question rather than a specified conjectural bound or algorithmic guarantee. The source does not specify a representation of \(P\), whether the dimension is fixed, an admissible computational cost, or an error criterion for estimates. Consequently, no particular efficiency or approximation requirement can be included without further hypotheses.
\subsection{Short English statement}
Given a polytope \(P\subseteq\R^n\), develop methods to compute or estimate how many integer lattice points it contains, including points on its boundary. The supplied question leaves the input representation, computational requirements, and approximation accuracy unspecified.
\subsection{Sources}
\begin{itemize}
\item Supplied local UnsolvedMath registry record: permanent ID \(30001380\), title ``Counting Integer Points in Polytopes,'' source problem code \texttt{catalogue entry}, in the ``Oberwolfach Reports Open Problems'' collection. The record gives proposed year 2009 and identifies no individual proposer.
\item Supplied report citation: \emph{Convex Geometry and its Applications} (2010), Oberwolfach Reports, pp.~2851--2926, DOI \texttt{10.4171/owr/2009/53}; \url{https://doi.org/10.4171/owr/2009/53}. The background describes Oberwolfach Reports as documenting mathematical workshops and their open-problem discussions. The literature-source label for this reference is ``original Oberwolfach report 4136.''
\item Supplied literature reference, labeled ``Barvinok fixed-dimension counting'': \url{https://doi.org/10.1287/moor.19.4.769}.
\item Supplied literature reference, labeled ``Barvinok--Hartigan Gaussian approximation'': \url{https://arxiv.org/abs/0903.5223}.
\item Literature assessment dated exactly \texttt{2026-08-21}: the supplied curation labels the problem ``partially solved.'' It reports general generating-function methods and efficient algorithms in fixed dimension, Gaussian and entropy approximations for broad asymptotic families, and no single unrestricted efficient method answering every instance. This is supplied dated curation derived from the OWR record, not an independently verified assessment; its time-sensitive status and cited claims require independent verification.
\item The supplied statement-status label is \texttt{corrected\_verified}; its accompanying explanation says that the polytope and lattice-point count were expressed in standard mathematical notation without changing the question. This records supplied editorial metadata, not independent mathematical or source verification.
\end{itemize}
\end{document}
